% Transcription — fonds Alexandre Grothendieck, shelfmark 57, batch 2 (pages 21-40).
% Established from the facsimile archives/batches/57/batch-02.pdf, per the skill
% transcribe-grothendieck.
%
% Pass: Opus 5.5 (claude-opus-5-5), 2026-10-03 — first pass, unchecked against the pages by a human.
%
% Fifteen of the twenty pages are transcribed. Page 23 is a cover sheet in his
% hand (« Rigidité », « Groupes V »), which becomes the section heading and gets
% no \page{}. Pages 35, 37 and 39 are unrelated typed versos (a stationery price
% list) and are absent.
%
%   Pages 21-22 — end of his typed letter to Murre begun on pp. 19-20
%     (separatedness of Corr, pre-ét-representability of Pic, sending notes).
%   Pages 24-34, 36, 38, 40 — manuscript « Rigidité et groupes alg. », his own
%     pagination 1-3 on pp. 24, 26, 28: rigidity theorem for morphisms
%     X_1 x_S X_2 -> G, lemmas 1-2, corollaries (abelian schemes commutative,
%     unique group law), Picard-functor reformulation (Corr separated),
%     Théorèmes 1-2 (p. 32, struck through), lemma 1 on varieties with a
%     multiplication, lemma 3 (deformation of the group law over an artinian
%     base, obstruction in H^1), and the final theorem on p. 40.
%
% Notation fixed for the batch: the underlined O of the structure sheaf is
% \underline{O}; underlined Pic/Corr are \underline{\mathrm{Pic}}, \underline{\mathrm{Corr}};
% the inverse map is written with a superscript \vee as he does.

\documentclass[11pt,a4paper]{article}
\input{../preamble/grothendieck}

\folder{57}
\batch{2}
\pages{21}{40}
\foldertitle{Picard : tapuscrit annoté (s.d.), lettres (s.d., 1962).}
\dating{1962-[vers 1968]}
\watermark{Édition de démonstration}

\begin{document}

\page{21}
\note{Fin d'une lettre dactylographiée en anglais, commencée avant ce lot (« Dear Murre », p.~19 du fonds) : c'est sa lettre à J.~P. Murre, reprise ici à partir de la p.~21. La flèche que la machine ne portait pas est restituée comme ajout de l'éditeur.}
simplicity the fibers of $X$ géoémetrically normal, and hence $X$ normal !), hence $D$ is linearly equivalent to $0$, what we wanted to prove.

I am convinced however that $\underline{\mathrm{Corr}}$ is always separated (with the usual assumptions of properness, flatness, and direct image of the structure sheaf, for both factors, of course). This is easily seen to be true if the $\underline{\mathrm{Pic}}$ functor of either factor is representable, by a simple use of dimension theory (namely, we have a morphisme $X$ \supplied{$\to$} $\underline{\mathrm{Pic}}_{Y/S}$ whose image has a general fiber of dimension zero, hence the same holds for the special fiber …). But it is true also, by an immediate adaptation of the same argument, if we suppose only that $\underline{\mathrm{Pic}}_{Y/S}$ is pre-ét-représentable say, i.e.\ is a quotient of a representable functor $Q$ by an étale equivalence relation (in fact, quasi-finite and flat would do as well), with $Q$ locally of finite type over $S$. Now this assumption is certainly satisfied if $Y$ is \emph{projective} over $S$, as one sees by using the reprsentation of $\underline{\mathrm{Pic}}_{Y/S}$ (or rather big open pieces of it) as the quotient of a suitable scheme of immersions of $Y$ into some $P^r$, by the action of the projective group operating freely, and taking a quasi-section of the corresponding equivalence relation … On the other hand, if one does not assume $X$ nor $Y$ projective over $S$, one may think of using Chow's lemma; as $S$ is the spectrum of a discrete valuation ring, one does not loose flatness in using Chow's lemma, unfortunately one will loose however, I am affraid, the assumption $H^0(X_0, \underline{O}_{X_0}) \simeq k(s)$, and I am affraid that this will make serious technical trouble. Another interesting approach, via topology, is to try to prove that under the usual assumptions on $X$, the "specialization morphisme" from the fundamental group of the general geometric fiber to the one of the special fiber has an image of finite index -- or at least that this is so after

\page{22}
making the groups abelian. It seems to me that the latter statement can be proved via the Picard functor, when $X$ is assumed projective over $S$.

I am sending you some notes, including a sketch of the proof of the theorem of representability of unramified functors, although I do not think the latter can be of any use to you, as i have a hard time myself to read them. I think the notes you took when we discussed the matter a few month ago should be much more detailed; anyhow, there are certainly no simplifications in my notes relative to yours, the inverse is more plausible.

Sincerely yours

\section{Rigidité}

\note{Titre de sa main, sur la chemise de la p.~23, qui porte « Rigidité » et au-dessous « \uncertain{Groupes} V », le V souligné  ; la chemise n'a pas de \emph{page} propre ici.}

\page{24}
\emph{Rigidité et groupes alg.} 1\note{titre en capitales en tête de page ; « ET » est écrit sur un mot biffé illisible, et le « 1 » à droite est la pagination de l'auteur.}

\marginal{« Théorème \ill{} »\note{note encadrée, écrite en oblique dans la marge supérieure gauche, reliée par un trait au mot « Théorème » ; au-dessous, un petit schéma : $X_1 \times_S X_2$ avec ses projections $p_1, p_2$ vers $X_1, X_2$, et $\varphi_1, \varphi_2$ vers $S$.}}
\emph{Théorème.} Soient $S$ un préschéma, $\varphi_i \colon X_i \to S$ ($i = 1, 2$) deux préschémas \struck{loc.\ de type fini sur $S$,} de type fini sur $S$, \add{$\varepsilon$ une section de $X_1$ sur $S$,} \struck{$G$ \ill{} schéma en groupes} \struck{on suppose} soient $p_i$ ($i=1,2$) les deux projections de $X_1 \times_S X_2$ sur les facteurs. On suppose :

a) $p_1$ est propre, et $p_{1*}(\underline{O}_{X_1 \times X_2}) \simeq \underline{O}_{X_1}$.

\struck{b) $\varphi$ \ill{} \ill{} $U$ de $\varepsilon(S)$ \ill{} $X_1$ \ill{} $X_1 \times_S X_2$ \ill{}}

b) $p_1$ est universellement ouvert, \struck{\ill{}} $X_1$ plat sur $S$ \add{à fibres géom.\ \uncertain{primaires}}.

[\emph{Ex.} $\varphi_2$ est propre, plat et $s \in S \Rightarrow k(s) \simeq H^0(X_{2,s}, \underline{O}_{X_{2,s}})$ ; \struck{« primaires » i.e.\ sans cycles ass.\ immergés,} $\varphi_1$ \struck{plat} à fibres géom.\ irréductibles]. Alors pour tout préschéma en groupes $G$ sur $S$, séparé sur $S$, et tout $S$-\uncertain{morphisme}
\[
  f \colon X_1 \times_S X_2 \longrightarrow G
\]
on a
\[
  f = (f_1 p_1)(f_2 p_2)
\]
\note{le produit est d'abord écrit $(f_1 \cdot p_1^{*})$ puis retouché ; la lecture $(f_1 p_1)(f_2 p_2)$ est celle que la suite confirme.}
où
\[
  f_1 \colon X_1 \to G, \qquad f_2 \colon X_2 \to G
\]
\note{il écrit « $f_2 \colon X_2 \to X$ », par lapsus pour $G$.}
sont des $S$-homomorphismes\note{sic : il faut entendre « $S$-morphismes ».}, déterminés de façon unique par la condition supplémentaire

$f_1 \varepsilon_1 = $ section unité \struck{pour \ill{}}, \add{et donc} $f_2 = f \circ (\varepsilon_1 \times_S \mathrm{id}_{X_2})$ \struck{$f_2 = f \circ (\varepsilon_1 \times_S \mathrm{id}_{X_2})$} \struck{[ \ill{} $\varphi_2$ soit une \ill{} $f_2 = $ ]} \uncertain{de plus} \ill{} $X_2$ \ill{}.

\marginal{\uncertain{Hypothèses} \ill{} \ill{} par la \ill{} \uncertain{commutative} plus haut}
\emph{Corollaire.} Supposons \ill{} \uncertain{section} $\varepsilon_2$ sur $S$. Alors $f$ se met de façon unique sous la forme
\[
  f = f_1 p_1 + f_2 p_2 + \sigma \varphi
\]
\uncertain{(tel que $f_i \varepsilon_i = $ section unité)}, où $f_i \colon X_i \to G$ sont des $S$-morphismes, $\sigma \colon S \to G$ une section de $G$ sur $S$, [$\varphi$ désignant \struck{$\varphi_1 \times_S \varphi_2$} la projection $X_1 \times_S X_2 \to S$].

\page{25}
\emph{Démonstration du th.}

Soit
\[
  f_2 = f \circ (\varepsilon_1 \times_S \mathrm{id}_{X_2})
\]
et considérons
\[
  f' = f - f_2 p_2 \colon X_1 \times_S X_2 \longrightarrow G ;
\]
\marginal{petit diagramme : $X_1 \times_S X_2 \to X_1$ et $X_1 \times_S X_2 \to X_1 \times G$, $X_1 \times G \to X_1$.}
\struck{\ill{} \ill{} \ill{} a) et}
\ill{} \ill{} \ill{} \ill{} \ill{} \ill{} $f' \circ (\varepsilon_1 \times \mathrm{id}_{X_2}) = $ morphisme unité $X_2 \to G$, i.e.\ pour tout $x_1 \in \varepsilon_1(S)$, le morphisme induit
\[
  f' \colon k(x_1) \otimes_{k(s)} X_2 \longrightarrow G
\]
\ill{} \ill{} \ill{} \ill{} \uncertain{constant}. En vertu de l'hyp.\ a) et du \uncertain{lemme} \ill{} ci-dessous, il en résulte qu'il existe un voisinage $U_1$ de $\varepsilon_1(S)$ tel que $f' \mid U_1 \times_S X_2$ soit de la forme $f_1 p_1$, où $f_1 \colon U_1 \to G$ est un $S$-morphisme \uncertain{convenable}, \ill{} \ill{} pour $x_1 \in U_1$, $f' \circ$ \ill{} \ill{} $k(x_1) \otimes_S X_2$ \ill{} \ill{} \ill{} \ill{} \ill{} $G$. En vertu du lemme~2 \struck{\ill{}} \add{et $p_1$ étant universellement \uncertain{ouvert}} b), \ill{} \ill{} \ill{} \ill{} \ill{} ci-dessus, \uncertain{cela} \ill{} \ill{} \ill{} \ill{} $x \in \overline{U}_1$, i.e.\ $\forall x \in X_1 \times_S X_2$ \ill{} \ill{} c). En vertu du lemme~1, on \ill{} \ill{} \ill{} prendre $U_1 = X_1$.

Cela prouve l'existence de $f_1$, $f_2$ ayant les \uncertain{propriétés} ($\ast$), l'unicité \uncertain{pour} \ill{} \ill{} \struck{[ $f_1 \varepsilon_1 = $ section \ill{} \ill{} \ill{} ] \ill{} de la \ill{} de $f_2 \varepsilon$, \ill{} \ill{}} \struck{\ill{} \ill{} \ill{} \ill{} \ill{}} \struck{[\ill{} $f_1$ \ill{} \ill{}]} \ill{} \ill{} \ill{} \ill{}.

\struck{$f = f_1 p_1^{*} f_2 p_2$, \ill{} $f_1 \varepsilon_1$ \ill{} $f_2$ \ill{} \ill{} \ill{} \ill{} \ill{} $\varepsilon_1 \times \mathrm{id}_{X_2}$}
\struck{\ill{} \ill{} \ill{} \ill{}}
\add{\uncertain{morphisme} \uncertain{unité}} $(f_1 \varepsilon_1)\varphi_2 = $ \struck{section} \ill{} \ill{} \ill{} \ill{}. D'ailleurs, \uncertain{dans} \ill{} \ill{} \ill{} \ill{} \ill{} \ill{} \ill{} \ill{} \ill{} \ill{} \ill{} \ill{}
\[
  f_2 = f \circ (\varepsilon_1 \times \mathrm{id}_{X_2})
\]
\note{le bas de la page, très surchargé de ratures et d'ajouts, n'est lu que par fragments.}

\page{26}
\note{En haut à droite, « 2 » : pagination de l'auteur.}

\emph{Lemme 1.} Soient $f \colon X \to Y$ un morphisme propre, $Y$ loc.\ noeth., tel que $f_*(\underline{O}_X) \simeq \underline{O}_Y$, \struck{et soit} \add{\uncertain{soit}} $h \colon Z \to Y$\note{un symbole biffé avant $Y$.} un morphisme séparé loc.\ de type fini, \struck{$Z$ préschéma de type fini,} $g \colon X \to Z$ un $Y$-morphisme, \struck{\ill{}} $y \in Y$, \uncertain{supposons} que $g(X_y) \subset Z_y$ soit fini. Alors il existe un \add{\uncertain{ouvert}} voisinage $U$ de $y$ et une section $\sigma$ de $Z$ sur $U$, telles que
\[
  g \mid f^{-1}(U) = \sigma \circ (f \mid f^{-1}(U)) .
\]
Pour $U$ donné, $\sigma$ est unique.

\marginal{\ill{} \uncertain{important} \ill{} sous forme \uncertain{corollaire} plus bas}

\emph{Dém.} Considérons $g(X) = X' \subset Z$. Il est propre sur $Y$, et sa fibre $X'_y$ en $y$ est finie, donc $\exists\, U \ni y$ tq $X' \mid U$ soit fini sur $U$. On peut supposer $U = Y$. Soit $\mathcal{A} = g_*(\underline{O}_X)$, c'est un faisceau cohérent d'algèbres sur $Z$, \uncertain{définissant} un $Z'$ fini sur $Z$, et on a une factorisation canonique $g = g'' g'$. On a\note{au-dessus de la ligne : « $\mathcal{A} = g_*(\underline{O}_X)$ ».}
\[
  \underline{O}_Y \to h_*(\mathcal{A}) \subset f_*(\underline{O}_X) \simeq \underline{O}_Y ,
\]
\note{le premier terme est d'abord écrit $g''(\underline{O})$, biffé, puis $\underline{O}_Y$.}
donc
\[
  (h \circ g'')_*(\underline{O}_{Z'}) \simeq \underline{O}_Y ,
\]
i.e.\ $Z' \to Y$ est un \emph{isom.} On prend alors $\sigma = g''(h g'')^{-1}$. L'unicité de $\sigma$ est évidente. De cette unicité \ill{} \ill{}

\note{diagramme de la marge gauche :}

\begin{tikzcd}[column sep=small, row sep=small]
 & Z' \arrow[dr, "g''"] & \\
X \arrow[ur, "g'"] \arrow[rr, "g"] \arrow[d, "f"'] & & Z \arrow[dll, "h"] \\
Y & &
\end{tikzcd}

\emph{Corollaire.} Sous les conditions du lemme, l'ens.\ $U$ des $y \in Y$ satisfaisant la condition $g(X_y)$ fini est ouvert, et il existe une section unique

\page{27}
$\sigma$ de $Z$ sur $U$ telle que $g \mid f^{-1}(U) = \sigma (f \mid f^{-1}(U))$.

\emph{Lemme 2.} Soient $f \colon X \to Y$ un morphisme de type fini \emph{universellement ouvert}, \struck{dans l'\ill{}} $h \colon Z \to Y$ un morphisme \add{loc.} de type fini, \struck{\ill{} \ill{} \ill{}} $g \colon X \to Z$ un $Y$-morphisme, alors l'ens.\ $U$ des $y \in Y$ tels que $g(X_y)$ soit fini est \uncertain{ouvert} \struck{\ill{}} \emph{fermé}.

\uncertain{En effet} \uncertain{on se ramène} \ill{} \ill{} $Y$ \uncertain{noethérien}, \struck{\ill{}} $X \to Z$ \uncertain{dominant}, \ill{} \uncertain{constructibles} \ill{} \ill{} \ill{} \ill{} \ill{} $Y$ \ill{} le \uncertain{spectre} \ill{} \ill{} \ill{} \ill{} \uncertain{discrète}, et \ill{} \ill{} \ill{} \ill{} \ill{} $y$ est \ill{} $U$. On peut \ill{} \ill{} \uncertain{supposer} $X$ \add{\uncertain{il}} irréductible, \struck{donc} \uncertain{dominant} \ill{} \ill{} \uncertain{réduit} $Y$, donc $\overline{g(X)}$ est irréductible et \uncertain{domine} $Y$, \struck{\ill{}} la fibre générique \ill{} \ill{} \ill{} \ill{} la fibre spéciale, \ill{} \ill{} \ill{}. \uncertain{cqfd}.

\emph{Corollaire 1} \struck{\uncertain{Supposons} $p_*(\underline{O}_{X_1})$}\note{biffure ondulée.} Dans la conclusion du \uncertain{th.}\ il suffit de \struck{\ill{} \ill{} \ill{}} \add{\uncertain{supposer} \ill{}} \ill{}

\emph{Remarque.} \struck{Le lemme 2 est \ill{}} \ill{} \ill{} \struck{\ill{} \ill{} \ill{} \ill{} \ill{} corollaire \ill{}.} $p \colon X_1 \times_S X_2 \to S$ est propre et $p_*(\underline{O}_{X_1 \times X_2}) \simeq \underline{O}_S$. \struck{\ill{} \ill{} \ill{} \ill{} \ill{} \ill{}} En effet, \ill{}
\[
  \text{\struck{$f_2 \; f = f \circ \varepsilon_1 p_1 - f \circ \varepsilon_2$}} \qquad f' = f - f_1 p_1 - f_2 p_2 \quad (f_i = f \circ u_i \ldots),
\]
il vient \uncertain{donc} \uncertain{pour} \ill{} \ill{} que $f'(\text{fibres})$ est fini, \uncertain{donc} \ill{} \ill{} \ill{} \ill{} \ill{} \ill{} \ill{} \ill{} \ill{} \uncertain{et appliquer} lemme~2.

\page{28}
\note{En haut à droite, « 3 » : pagination de l'auteur.}

\emph{Corollaire 2.} Soit \struck{$A$} $A$ un schéma abélien sur $S$ loc.\ noethérien, et $G$ un préschéma en groupes loc.\ de type fini sur $S$. Alors \struck{\ill{}} \struck{\ill{}} tout homomorphisme \add{de $S$-préschémas} de $A$ dans $G$ est de la forme (de façon unique)
\[
  f = u + c
\]
où $u$ est un homomorphisme de groupes, et où $c$ est un homomorphisme « \uncertain{constant} » [i.e.\ \uncertain{produit} d'une section de $G$ sur $S$ \ill{}]. On a nécessairement
\[
  c = f \circ \varepsilon_A ,
\]
où $\varepsilon_A$ est la section unité de $A$.

\marginal{\ill{} \ill{} \ill{} \uncertain{commutatif}}
\uncertain{Donc} \ill{} si $f \circ \varepsilon_A = \varepsilon_G$, $f$ est un homomorphisme de schémas en groupes.

\emph{Corollaire 3.} Soit $A$ un schéma abélien sur $S$. Alors $A$ est commutatif. Toute loi de composition sur $A$, telle que $\varepsilon_A$ \struck{\ill{}} soit \uncertain{élément} unité \uncertain{bilatère}, est la loi de composition donnée de $A$.

Il suffit de prouver cette dernière \uncertain{assertion}. Or si $f \colon A \times A \to A$ est une loi de composition, en vertu du cor.~1 :
\[
  f = [f \circ (\varepsilon \times \mathrm{id}_A)]\, p_1 + [f \circ (\mathrm{id}_A \times \varepsilon)]\, p_2 + [f \circ (\varepsilon_1 \times \varepsilon_2)]\, p\,,
\]
\note{entre le deuxième et le troisième terme, un terme biffé illisible.}
i.e.\ $f = p_1 + p_2$, c'est la loi de composition donnée. OK.

\page{29}
\emph{Reformulation du th.}

\emph{Th.} $X$ propre et plat sur $S$, $H^0(X_s, \underline{O}_{X_s}) \simeq k(s)$ pour tout $s \in S$, ($S$ loc.\ noeth.), $\varepsilon_X$ une section de $X/S$, \struck{$\mathrm{Hom}(X,$} $G$ un foncteur sur \struck{\ill{}} $(\mathit{Sch})_{/S}^{\circ}$ à valeurs dans la catégorie des groupes, $f \colon X \to G$ un $S$-homom.\ \uncertain{fonctoriel}. \struck{\ill{} \ill{} \ill{} \ill{} \ill{}} \struck{\ill{} \ill{} on suppose donné}

(i) Si $S_0$ est un sous-préschéma \emph{\uncertain{réduit}} de $S$ ayant même ensemble sous-jacent, et si \struck{$G$} l'homom.\ $f_0 \colon X_0 \to G_0$ est trivial, alors $f$ l'est [pourvu que $G$ \ill{} \ill{} \ill{} \uncertain{petite catégorie})].

(ii) Si de plus $G$ commute aux limites projectives \add{[cas de la localisation]} d'anneaux noethériens, \ill{} \ill{}, aux \uncertain{limites} d'\uncertain{anneaux} \uncertain{complets}\note{passage très surchargé.} \ill{} \ill{} \ill{} aux morphismes $\mathrm{Spec}(\widehat{A}) \to \mathrm{Spec}(A)$ des morphismes \ill{} \ill{} $G$-\ill{}, \uncertain{alors} si $f_s = f \times_S k(s)$ est triviale, $f$ est triviale dans un voisinage de $s$.

(iii) Si $G$ est représentable \add{\uncertain{séparé}} sur $S$, alors sous les conditions de (ii), \struck{\ill{}} $f$ est trivial sur la composante connexe de $s$.

\uncertain{Formulation} \ill{} \uncertain{lemme}, $\underline{\mathrm{Hom}}_S^{\bullet}(X, G)$ est \emph{non ramifié} sur $S$.

\page{30}
\emph{Th.} Soient $f \colon X \to S$ propre et plat, $\underline{O}_S \simeq f_*(\underline{O}_X)$ universellement, $S$ loc.\ noeth., $G$ foncteur $(\mathit{Sch})_{/S}^{\circ} \to (\mathit{Ens})$\note{ici « $(\mathit{Ens})$ » ; à la p.~29 il demandait des valeurs dans les groupes.} :

a.\ $G$ formellement représentable \ldots

\struck{b) $G$ « compatible avec limites d'anneaux \uncertain{artiniens} » \ill{}}\note{le premier b) est biffé ; au-dessus, entouré, « \uncertain{de notre} \ill{} \ill{} ».}

b) $G$ compatible avec \uncertain{passage aux complétés}\note{lecture très incertaine de la fin de la ligne.} \ill{} \ill{} \ill{} \uncertain{plats} et \uncertain{qui} \uncertain{respectent}

c) $G$ compatible avec \emph{limites} d'\add{anneaux} \struck{\ill{} \ill{}}

\marginal{(Par exemple, il \uncertain{suffit} que $G$ soit représentable \ill{} et \ill{}\note{note encadrée, en haut à droite, au crayon : « Par exemple il suffit que $G/S$ soit \ill{} \ill{} \ill{} » ; elle est lue de façon trop incertaine pour être transcrite.}}
\marginal{Conditions \uncertain{vérifiées} si $G$ est représentable loc.\ de type fini sur $S$\note{accolade en marge droite englobant a) à c), écrite verticalement.}}
\marginal{\uncertain{Utiliser} la \ill{} \ill{} \ill{}\note{marge gauche, en oblique.}}

$u \colon X \to G$, i.e.\ $u \in G(X)$. Soit $U$ l'ens.\ des $s \in S$ tels que $u_s \in G(X_s)$ soit l'image d'un élément de $G(k(s))$, i.e.\ $u_s \colon X_s \to G$ se factorise en
\[
  X_s \to k(s) \to G .
\]
Alors

1) $U$ est ouvert [et \emph{fermé} \uncertain{si} \ill{} \ill{} \ldots \uncertain{(critère valuatif)}]\note{passage dans une autre encre, avec renvoi.}

2) Il existe un \struck{morphisme} $v \in G(U)$ unique tel que l'image de $v$ dans $G(X \mid U)$ soit $u \mid (X \mid U)$\note{il écrit $u \mid (X \mid \overset{\circ}{U})$ ; le petit signe sur $U$ est douteux.}, \struck{etc.}

\emph{Corollaire 1.} Si $S$ est connexe, et s'il existe un $s \in S$ tel que $u_s \colon X_s \to G$ se factorise par $k(s)$, alors $u$ se factorise par
\[
  X \xrightarrow{\ f\ } S \xrightarrow{\ v\ } G .
\]

\emph{Corollaire 2.} Supposons que $X$ et $G$ soient munis de sections sur $S$, et que $u$ soit compatible avec les sections unités. Si $S$ est connexe, et si pour \emph{un} $s \in S$, $u_s$ est trivial, alors $u$ est trivial.

\emph{Corollaire 3.} $X$ sur $S$, \struck{\ill{}} \add{\uncertain{du cor.~2}} $G$ sur $S$ comme \struck{\ill{}}, \struck{\ill{} connexe à sections,} $Y$ sur $S$ loc.\ de type fini \add{$Y$ connexe}, \struck{sur} $S$, et $u \colon X \times_S Y \to G$ un homom.\ \uncertain{fonctoriel} \struck{\uncertain{fidèlement} \ill{}} \ill{} l'homom.\ trivial sur $\varepsilon_X \times_S Y$, et s'il existe $y \in Y$ tel que $u_y \colon X \times k(y) \to G$ soit trivial, alors $u$ est trivial.

\page{31}
\emph{Cor.~4.} Sous les conditions du cor.~2, soit de plus $Y$ loc.\ de type fini sur $S$, \add{à fibres géom.\ connexes, $Y$} muni d'une section $\varepsilon_Y$. Soit $u \colon X \times Y \to G$ induisant les morphismes triviaux $\varepsilon_X \times Y \to G$ et $X \times \varepsilon_Y \to G$, alors $u$ est trivial.

\emph{Cor.~5} \add{(\uncertain{Weil})} Sous les conditions du cor.~4, supposons \struck{\uncertain{foncteur} en groupes, \uncertain{les sections}} de plus $G$ un \ill{}. Alors tout morphisme $u \colon X \times_S Y \to G$ se met de façon unique sous la forme
\[
  u = u_1\, \mathrm{pr}_1 + u_2\, \mathrm{pr}_2 + v\, p
\]
\note{le dernier terme est surchargé.}
pris dans le groupe $X \times_S Y = \underline{X}_Y$\note{lecture incertaine de cette égalité.}, $p \colon X \times_S Y \to S$ le morphisme structural, $u_1 \colon X \to G$, $u_2 \colon Y \to G$\note{il écrit $u_2 \colon X \to G$.}, $v \colon S \to G$ des $S$-morphismes convenables.

[\uncertain{Demander} Mumford]\note{cette ligne et la suite de la page sont d'une autre encre.}

\emph{Question.} Soit $S$ (\uncertain{noethérien} \ill{} \ill{} \ill{} \ill{} \ill{} $S$ spectre d'un anneau de val.\ discrète \uncertain{géom.\ complet}), $X$, $Y$ propres et plats sur $S$, tels que $f_*(\underline{O}_X) \simeq \underline{O}_S$ et $g_*(\underline{O}_Y) \simeq \underline{O}_S$ universellement, munis de sections $e_X$ et $e_Y$. Soit $\underline{L}$ un Module inversible sur $X$, \struck{\ill{}} trivialisations données sur $e_X \times_S Y$ et $X \times_S e_Y$, [compatibles sur $e_X \times_S e_Y$], et supposons qu'il existe un ouvert non vide $U$ tel que $\underline{L} \mid X \times_S Y \mid U$ soit trivial. Alors $\underline{L}$ est trivial [En d'autres termes, $\underline{\mathrm{Corr}}_S(X, Y)$ est \uncertain{toujours} \emph{séparé sur $S$}] ; \uncertain{c'est vrai} si \struck{\ill{}} $\underline{\mathrm{Pic}}_{X \times_S Y}$ \ill{} \ill{} \ill{}, par exemple si $\underline{\mathrm{Pic}}_{X/S}$ et $\underline{\mathrm{Pic}}_{Y/S}$ \ill{} \ill{} \ill{}]. Rejoint l'affirmation si $\underline{\mathrm{Pic}}_{X/S}$ ou $\underline{\mathrm{Pic}}_{Y/S}$ existe sur les anneaux de val.\ discrète complets. \uncertain{Désirons} \struck{\ill{} \ill{} \ill{} \ill{} \ill{}}

\marginal{Il suffit que $\underline{\mathrm{Pic}}_{X/S}$ ou $\underline{\mathrm{Pic}}_{Y/S}$ soit pré-ét-représentable, \uncertain{en particulier} \uncertain{suffit} que $X$ ou $Y$ soit projectif sur $S$.}
\marginal{\uncertain{Si} $f$ \uncertain{ou} $g$ \ill{} \ill{} \ill{} \uncertain{Picard} \ill{} sur $S$.\note{deux notes de la marge inférieure gauche, reliées par des renvois au texte ; cf.\ la lettre des pp.~21--22, qui énonce les mêmes conditions.}}

\page{32}
\emph{Théorème 1} \add{(Définition des schémas abéliens)}. Soient \struck{$S$} $S$ préschéma loc.\ noethérien \emph{normal}\note{lecture incertaine du dernier mot, souligné.}, $X$ $S$-préschéma \emph{propre} et \emph{simple}, $e$ une section de $X$ sur $S$, \uncertain{supposons} que pour tout $s \in S$ tel que $X_s = f^{-1}(s)$ soit un schéma abélien sur $k(s)$. Alors il existe \struck{\ill{}} sur $X$ une structure de schéma abélien sur $S$, admettant $e$ comme section unité\note{la fin de la phrase est en partie biffée.}.

\emph{Cor.} \uncertain{Même} si $S$ \uncertain{n'est} pas normal, une structure abélienne est \uncertain{déterminée} par la \uncertain{connaissance} de \struck{la} section unité. (\uncertain{Cela} \ill{} \ill{} \ill{})

\note{À partir d'ici et jusqu'au bas de la page, le texte est barré d'un réseau de traits obliques ; il se lit sous la rature.}

\emph{Théorème 2.} Soient $S$ \struck{\ill{}} \uncertain{noethérien} \add{\uncertain{de type fini}}, $X_1$ et $X_2$ \add{\uncertain{deux}} \uncertain{préschémas} \ill{} sur $S$ ($\varphi_i \colon X_i \to S$). On suppose :

(i) les $\varphi_i$ plats et $\varphi_{2*}(\underline{O}_{X_2}) = \underline{O}_S$, \emph{ou} $\varphi_2$ séparable ;

(ii) les fibres de $X_1$ \ill{}, et $X_1$ est connexe.

$X_1$ \uncertain{muni} d'une section \uncertain{$s_1$}\note{« $s_1$ » cerclé.}.

\struck{Alors \ill{}} \uncertain{Soit} $G$ schéma en groupes sur $S$, \struck{abélien}, \add{\uncertain{considérons}}
\[
  u \colon X_1 \times_S X_2 \to G ;
\]
$\exists\, u_1 \colon X_1 \to S$\note{sic, pour $X_1 \to G$.} \uncertain{unique}, $u_2 \colon X_2 \to G$, \uncertain{tels que} \struck{\ill{}}
\[
  u = u_1\, \mathrm{pr}_1 + u_2\, \mathrm{pr}_2 .
\]

\emph{Cor.} \uncertain{Soit} \struck{$S$} $S$, \ill{} \struck{$X_1 = X_2 = $} $X$ \uncertain{un} \uncertain{schéma} \ill{} \uncertain{sur} $S$ \uncertain{à} fibres \uncertain{connexes}, \uncertain{et} \ill{}, $X/S$ plat, simple, $u \colon X \to G$ \uncertain{tel que} $u(e) = e$, alors $u$ \uncertain{est} \uncertain{un} homom.\ \uncertain{de groupes}.

\begin{tikzcd}
X \times X \arrow[r] \arrow[d, "{u \times u}"'] & X \arrow[d, "u"] \\
G \times G \arrow[r] & G
\end{tikzcd}

\page{33}
\note{Page presque vierge : trois lignes de programme, en tête.}

\emph{Unicité} (\struck{\ill{}}) \uncertain{$f$} \uncertain{comp.}\ \uncertain{Rem.}\ \uncertain{sur} \ill{} \ill{} \add{\uncertain{et commutativité}}

\emph{Schémas}

\page{34}
\emph{Lemme 1.} Soient $G$ \struck{\ill{}, \ill{} \ill{} connexe} \add{irréd.\ et \emph{propre}} sur $k$, $e$ point rationnel de $G$, $G \times G \to G$ loi de multiplication \struck{\uncertain{ayant}} \add{\uncertain{admettant}} $e$ comme unité, $G \xrightarrow{\ \vee\ } G$ \add{\uncertain{application}} \struck{\ill{} \ill{} à droite} telle que
\[
  x y^{\vee} = e \iff x = y .
\]
\emph{Alors} \uncertain{c'est} la loi de composition est associative et commutative [donc $G$ est une variété abélienne].

\emph{Démonstration.} Considérons,
\[
  \text{\struck{$f(a,b,c) = [(ab)c]\,[a(bc)]^{\vee}$}}
\]
\[
  f(a,b) = (ab)(ba)^{\vee} .
\]
On a
\[
  f(e, x) = f(x, e) = e ,
\]
donc
\[
  \begin{cases}
    f(a, x) = f(e, x) = e & \text{pour } a \in U \text{ voisinage de } e, \\
    f(x, a) = f(x, e) = e & \text{pour } a \in U',
  \end{cases}
\]
\note{une accolade et une flèche renvoient ces deux lignes à « rigidité ».}
\struck{donc $f(x,x) = f(e,x) = f(x,e) = f(e,e) =$}

donc par \emph{\uncertain{prolongement}} [\ill{} \ill{} \ill{}] on a $f(x, y) = e$ pour tout $x, y$, d'\uncertain{où} \uncertain{commutativité}.

Donc $(ab)(ba)^{\vee} = e$ i.e.\ $ba = ab$.

Pour l'associativité, on fait
\[
  f(a,b,c) = [(ab)c]\,[a(bc)]^{\vee}
\]
\note{il écrit les deux crochets avec des virgules : « $[(ab, c]\,[a(bc,]^{\vee}$ ».}
\uncertain{on a} $f(a,b,c) = e$ si $a [\ldots c] = e$ \struck{\ill{} \ill{} \ill{}} \add{pour \ill{} \ill{} $a \in U$ \uncertain{voisinage} de $e$}

donc $f(x, y, z) = e$\note{lecture incertaine des arguments.}

\uncertain{lemme} $f(abc) =$\note{la page s'arrête sur ce signe d'égalité.}

\page{36}
\emph{Lemme 3.} Soit $G$ schéma simple et propre sur $S$ \add{$= \mathrm{Spec}(A)$} \struck{loc.\ noeth.}, \add{$A$} \uncertain{artinien} local, \struck{\uncertain{anneau}} $G$ muni d'une section $e$. \struck{Soit $S_0 = \mathrm{Spec}(A/\mathfrak{J})$,} \struck{\ill{} \ill{} \ill{} soit \ill{} \ill{}} \struck{\ill{} \uncertain{idéal} \uncertain{tel que} $\mathfrak{J}^2 = 0$} \add{\ill{} \ill{} $S_0$}\note{début très surchargé : une première rédaction, biffée, introduisait $S_0 = \mathrm{Spec}(A/\mathfrak{J})$ avec un idéal $\mathfrak{J}$ de carré nul ; la lecture définitive est incertaine.} \ill{} \ill{} \ill{} \uncertain{supposons} que
\[
  G_0 = G \times_S S_0
\]
soit un schéma abélien ayant $e_0$ comme section unité. \struck{Alors il \ill{}} \struck{Soit} Alors il existe une structure multiplicative unique sur $G$, \uncertain{induisant} celle de $G_0$ et \uncertain{pour} \uncertain{laquelle} $e$ \struck{\ill{} \ill{} \ill{} \ill{}} \add{\uncertain{soit} \ill{} \ill{}} (\uncertain{sic}) \ill{} \ill{} \uncertain{donc} $G$ est un schéma abélien par cette structure, de section unité $e$.

\emph{L'obstruction} \struck{\ill{}} \uncertain{à} \uncertain{construire} \uncertain{une} \ill{} \uncertain{supposons} $S$ \struck{\uncertain{noeth.}} \uncertain{affine}\note{lecture incertaine de cette ligne.}. L'obstruction \ill{} \uncertain{à un} prolongement de la multiplication $G_0 \times G_0 \xrightarrow{\ \pi_0\ } G_0$ est dans
\[
  H^1\bigl(G_0 \times G_0,\ \pi_0^{*}(\underline{\mathcal{V}}_{G_0/S_0}) \otimes \ill{}\bigr) = 
\]
\struck{\uncertain{par} \uncertain{Künneth} \ill{}}
\[
  H^1(G_0 \times G_0, \underline{O}_{G_0 \times G_0}) \otimes_{k} (\mathcal{V} \otimes_k \mathfrak{m}) .
\]
\note{il écrit d'abord $\otimes_{K_0}$, corrigé en $\otimes_k$ ; « $\mathfrak{m}$ » est une lecture probable de la dernière lettre.}
\uncertain{Künneth} \uncertain{s'applique}, il suffit de montrer que la restriction de \uncertain{l'obstruction} à $e \times G_0$ et $G_0 \times e$ est nulle. \uncertain{Or} c'est \uncertain{vrai}, car on \uncertain{peut} \struck{\ill{}} \add{\uncertain{relever}}, \uncertain{à} $H^1(G_0, \underline{O}_{G_0}) \otimes_k (\mathcal{V} \otimes_k \mathfrak{m})$, \ill{} \uncertain{les} \uncertain{obstructions} \uncertain{à} \uncertain{prolonger} $g \mapsto g\,g$ et $g \mapsto g\,e$, \uncertain{qui} \uncertain{sont} \uncertain{les} \uncertain{identités}, \uncertain{donc} \uncertain{se} \uncertain{prolongent}. \uncertain{Pour} \uncertain{l'unicité}, \ill{} \ill{} \ill{} \ill{} \uncertain{dans} $H^0(G_0 \times G_0, \underline{O}_{G_0 \times G_0}) \otimes_k (\mathcal{V} \otimes_k \mathfrak{m})$.

\marginal{\uncertain{Lemme 2} \ill{}\note{note de la marge gauche, en face du dernier alinéa.}}
\uncertain{Plus généralement}, \uncertain{si} \ill{} \ill{} \ill{} $X \to G$, \uncertain{l'obstruction} \uncertain{à} \uncertain{prolonger} est dans $H^1(X_0, \underline{O}_{X_0}) \otimes \mathcal{V} \otimes_k \mathfrak{m}$, et l'\uncertain{indétermination} \uncertain{dans} $H^0(X_0, \underline{O}_{X_0}) \otimes_k \mathcal{V} \otimes_k \mathfrak{m}$\note{les deux dernières lignes, au bas de la page, sont très serrées.}.

\page{38}
\marginal{\uncertain{du lemme 2}}
\emph{Cor.} Si \uncertain{dans} $H^0(X_0, \underline{O}_{X_0}) = k$, \uncertain{cette} \uncertain{indétermination} \uncertain{ne} \uncertain{dépend} \uncertain{que} de $k$ : elle \uncertain{est} \uncertain{donnée} \uncertain{par} \struck{\ill{}} \uncertain{l'image} \add{\uncertain{par} \uncertain{les}} \add{\uncertain{sections}} \uncertain{de} $X$ \uncertain{dans} \uncertain{les} \struck{\ill{} \ill{} \ill{} la} section \uncertain{unique} $e_0$ \struck{[\ill{} \ill{} \ill{} \uncertain{par} \ill{}, \ill{} \ill{} \ill{}]}. D'où \uncertain{une} \uncertain{caractérisation} \ill{} de l'unicité \ill{} \ill{} \uncertain{fixe} l'extension, \ill{} \ill{} la section \ill{} \ill{} \ill{} \ill{} $e$.

\uncertain{De plus}, \struck{\ill{} \ill{}} \add{\ill{} \uncertain{de}} \struck{\ill{} \ill{}} $G$ \uncertain{par} $e\,g$ et $g\,e$, \uncertain{elle} \uncertain{est} \uncertain{aussi} \uncertain{fixée} \uncertain{par} $e$, \struck{\uncertain{les} \ill{}} \uncertain{donc} $e$ \uncertain{est} \uncertain{donc} \uncertain{d'après} \uncertain{ce} \uncertain{qui} \uncertain{précède} :
\[
  e\,g = g, \qquad g\,e = g ;
\]
\uncertain{donc} $e$ \uncertain{est} \uncertain{section} \uncertain{unité} \uncertain{bilatère}. \uncertain{Par le} \uncertain{même} \uncertain{raisonnement}, \uncertain{on} \uncertain{voit} que $g\,g' = g'\,g$ et $(g\,g')\,g'' = g\,(g'\,g'')$, \uncertain{car} \uncertain{on} \uncertain{a} \uncertain{deux} \uncertain{morphismes} $G \times G \rightrightarrows G$ (\uncertain{resp.}\ $G \times G \times G \rightrightarrows G$) \uncertain{qui} \uncertain{coïncident} \uncertain{sur} \uncertain{le} \ill{}, \uncertain{et} \uncertain{transforment} $e \times e$ (resp.\ $e \times e \times e$) en $e$.

\emph{Cor.~1 du lemme 3.} Possibilité (\struck{\ill{}} \add{et unicité}) de l'extension \uncertain{d'une} \uncertain{structure} abélienne \uncertain{donnée} sur $G_0$, \uncertain{et} la section unité \ill{} [il \uncertain{suffit} \uncertain{d'exiger} $e \times e \to e$], \uncertain{lorsqu'on} \uncertain{se} \uncertain{donne} une \emph{section}.

\emph{Cor.~2.} \uncertain{Dans} \uncertain{le} \uncertain{cas} \uncertain{où} \uncertain{on} \uncertain{se} \uncertain{place} \uncertain{sur} \uncertain{un} \uncertain{anneau} local \uncertain{complet}.\note{la page s'arrête ici ; cette dernière phrase est lue avec beaucoup d'incertitude.}

\page{40}
\emph{Cor.~3.} Soit un schéma abélien \uncertain{sur} un $S$ loc.\ noeth., il y a une \emph{seule} structure multiplicative possible admettant la section $e$ comme section unité bilatère.

[\uncertain{S'obtient} aussi par \uncertain{d'autres} \uncertain{voies} \ill{} \ill{}]

À partir de l'unicité, on prouve l'existence \uncertain{quand} \uncertain{on est} \uncertain{sur} un \uncertain{anneau} \emph{local} noeth., \struck{$G$} simple et propre \uncertain{sur} $S$, à section \uncertain{unique}, $G_0$ schéma abélien sur $k$ : \uncertain{on} \uncertain{utilise} \uncertain{par} \uncertain{le} \uncertain{théorème} \emph{\uncertain{par}} \uncertain{de} \emph{descente}. Enfin, on \uncertain{doit} \uncertain{encore} \uncertain{prouver} \uncertain{de} \uncertain{l'énoncé} \uncertain{quand} $S$ \uncertain{est} loc.\ noeth.\ \uncertain{connexe}. \uncertain{Cela} \uncertain{est} \uncertain{conséquence} \uncertain{du} \uncertain{fait} \uncertain{que} l'\uncertain{ensemble} \uncertain{des} \uncertain{pts} \uncertain{où} \uncertain{la} \uncertain{fibre} \uncertain{est} \uncertain{abélienne} (\uncertain{qui} \uncertain{est} \uncertain{déjà} \emph{\uncertain{ouvert}}) \uncertain{est} \uncertain{aussi} \emph{fermé}. \uncertain{On} \uncertain{se} \uncertain{ramène} \uncertain{au} \uncertain{cas} \uncertain{où} $S = \mathrm{Spec}(V)$, $V$ \uncertain{anneau} \uncertain{de} \uncertain{valuation} discrète. Alors \uncertain{c'est} \uncertain{connu}. \struck{\ill{}} [\uncertain{Lemme} \uncertain{de} \uncertain{décalage}] D'où

\emph{Théorème.} Soit $S$ loc.\ noeth.\ connexe, $G$ simple et propre sur $S$ \uncertain{avec} \struck{section} \add{section} \uncertain{unique}. Si \emph{une} fibre est abélienne, il y a sur $G$ une structure mult.\ unique \uncertain{d.g.}\ \struck{\uncertain{$e \times e$}} \add{\uncertain{admettant} \uncertain{$e$} \uncertain{comme} \uncertain{unité}}, et elle est abélienne.

\note{L'argument s'arrête en bas de la p.~40 sur l'énoncé du théorème ; la suite, s'il y en a une, est au-delà de ce lot.}

\end{document}
